\chapter{Datos e IA como Capacidad Estratégica}\label{ch:data-and-ai-strategy}

% Alcance: los datos como activo estratégico (*data asset* / ciclo de datos / efectos de red), IA construir vs comprar vs API,
%   gobernanza de IA y riesgo de modelo, ESG como problema estratégico de datos (doble materialidad).
%   Capítulo NUEVO (cierra la brecha de estrategia de datos e IA + parte de la profundidad ESG).
% Operacionaliza VRIN (\cref{ch:advantage-and-disruption}) y los *moats* (\cref{ch:growth}) para el contexto de IA.
% Borrador conforme al estilo de la colección: prosa continua + set-offs example/admonition; ejemplos trabajados con la SaaS de referencia.
% Este capítulo es CUALITATIVO: no se definen nuevas ecuaciones aquí. Todas las ecuaciones se referencian con \cref de etiquetas existentes.

\section{Los Datos como Activo Estratégico}\label{sec:data-asset}
% complexity: 3

¿Qué hace que un conjunto de datos sea competitivamente valioso? La intuición de que «más datos siempre es mejor»
sobreestima el argumento: la mayoría de las colecciones de datos alcanza un tope en la precisión predictiva mucho antes
de ser suficientemente grande para constituir un *moat*. El fenómeno más escaso — un \emph{data flywheel}
(ciclo de datos) — es un bucle de retroalimentación autorreinforcionante en el que un producto atrae
usuarios, los usuarios generan datos, los datos entrenan mejores predicciones, las mejores predicciones
mejoran el producto y el producto mejorado atrae más usuarios. Cada ciclo amplía la brecha frente a un
competidor que debe iniciar el bucle desde cero. La lógica del valor de red de \cref{eq:metcalfe}
explica por qué esto se compone: a medida que la base de usuarios crece de manera lineal, el espacio
combinatorio de señales de comportamiento crece de forma cuadrática, produciendo mejoras algorítmicas
que ningún conjunto de datos estático puede replicar.

La prueba estratégica para determinar si un \emph{data asset} (activo de datos) sostiene realmente la
ventaja es el marco VRIN de \cref{sec:rbv-vrin}: ¿es \emph{Valioso} (¿mejora predicciones que los
clientes valoran?), \emph{Raro} (¿pueden los competidores adquirir un corpus equivalente?),
\emph{Inimitable} (¿dependen los datos de una trayectoria histórica o son contingentes al pasado?)
y \emph{No sustituible} (¿puede un recurso distinto entregar el mismo resultado de producto?). Para
la mayoría de los datos operativos internos — registros de servidor, tickets de soporte, actividad CRM
genérica — la respuesta a al menos tres de esas cuatro pruebas es no, y la ventaja en datos es
ilusoria. Un verdadero *moat* de datos requiere que las cuatro sean sí
\parencite{agrawal2018prediction}.

\begin{example}[Datos de referencia multi-tenant en SaaS y la prueba VRIN]
Una plataforma SaaS con \money{4000000} ARR recopila telemetría de uso anonimizada en todos los
tenants: tasas de adopción de funcionalidades, tiempos de completación de flujos de trabajo,
frecuencias de errores y patrones de profundidad de sesión. El equipo de producto propone una
funcionalidad de \emph{benchmarking} (referencia) que muestra a cada cliente cómo se compara su
uso con la cohorte de pares de su vertical industrial. Se aplica la prueba VRIN de manera explícita.

\textbf{¿Valioso?} Sí: las referencias reducen la incertidumbre de los clientes sobre si están
extrayendo el valor total del producto, lo que aumenta la retención. La mejora en retention eleva
CLV directamente a través de \cref{eq:roic}: si el churn mensual cae de \qty{0.65}{\percent} a
\qty{0.55}{\percent} y la tasa de descuento es \qty{11}{\percent}, el valor presente de la base de
clientes aumenta de manera material.

\textbf{¿Raro?} Sí: los datos de uso multi-tenant con esta granularidad requieren una base de
clientes activa y con consentimiento a escala. Un nuevo participante no puede sintetizarlos; debe
ganárselos año tras año.

\textbf{¿Inimitable?} Sí: el corpus incorpora contexto histórico — patrones estacionales,
comportamiento durante migraciones de versión, rutas de onboarding específicas por industria —
que tomó años acumular. Replicar el algoritmo es sencillo; replicar el conjunto de entrenamiento, no.

\textbf{¿No sustituible?} Sí: ningún conjunto de datos públicamente disponible entrega telemetría
de flujos de trabajo SaaS en tiempo real, a nivel de funcionalidad y por cuenta. Las referencias
de encuestas y los informes de analistas son demasiado gruesos y demasiado tardíos para servir
al mismo propósito.

Las cuatro pruebas se superan. Este activo de datos sostiene ROIC\,>\,WACC (\cref{eq:roic}) más
allá de la vida de cualquier funcionalidad individual, porque cada nueva cohorte de clientes enriquece
el corpus y refuerza la señal de referencia que diferencia al producto.
\end{example}

\begin{admonition}{Advertencia}
La ventaja en datos se degrada bajo la regulación de portabilidad de datos. El artículo~20 del RGPD
otorga a las personas el derecho a recibir sus datos personales en un formato portátil y a
transmitirlos a un proveedor competidor. La Ley de Mercados Digitales de la UE extiende una
obligación estructuralmente similar a los guardianes de acceso. Un *moat* de datos cuya materia
prima son datos personales, en lugar de señales de comportamiento agregadas, puede erosionarse
por los derechos de portabilidad más rápido de lo que un nuevo ciclo de producto puede renovar
la ventaja competitiva. Audite si su ciclo de datos sobrevive a un régimen de portabilidad antes
de incorporarlo a su plan estratégico.
\end{admonition}

El ciclo de datos y el modelo de efectos de red de \cref{eq:metcalfe} comparten la misma lógica
estructural: el valor se acumula de manera no lineal a medida que crece la base de usuarios, pero
solo cuando el dato marginal mejora genuinamente el resultado central. Ambos mecanismos se conectan
con el diferencial de ROIC sostenible de \cref{eq:economic-profit}: el activo de datos es valioso
únicamente en la medida en que genera una mejora de producto que los clientes pagarán por conservar.
\textcite{agrawal2018prediction} desarrollan la economía de la predicción como insumo de costo
decreciente y sus implicaciones para la ventaja competitiva.

\section{Construir vs Comprar vs API en IA}\label{sec:ai-build-buy}
% complexity: 4

Toda decisión de inversión en IA se reduce a una versión de la misma pregunta: ¿en qué nivel del
stack debe la empresa poseer la capacidad en lugar de adquirirla? La economía de la IA ha cambiado
drásticamente desde principios de la década de 2020: el costo de ejecutar inferencia en un modelo
de propósito general potente ha caído en órdenes de magnitud, convirtiendo la opción de API en el
valor predeterminado correcto para la mayoría de las tareas y elevando el listón para justificar
la construcción de capacidad propietaria. La decisión se enmarca mejor como un mapeo de dos
dimensiones: ¿qué tan propietarios son los datos de entrenamiento de la empresa? y ¿qué tan
específica es la aplicación para esa tarea?

El diagrama 2×2 de \cref{fig:ai-build-buy} organiza los cuatro cuadrantes. Cuando tanto los datos
como la tarea son genéricos, usar un envoltorio de API es lo correcto: la economía de la inferencia
de commodities domina y no existe ningún *moat* sostenible. Cuando la tarea es genérica pero la
empresa posee datos propietarios, el ajuste fino (\emph{fine-tuning}) de un modelo base sobre esos
datos puede ofrecer una diferenciación moderada a un costo razonable. Cuando ni los datos ni la
tarea son propietarios, orquestar una capa de SaaS de proveedor suele ser la respuesta correcta:
la empresa captura la ganancia de productividad sin la carga de mantenimiento. Solo en el cuadrante
superior derecho — alta propietariedad de datos, alta especificidad de tarea — una construcción
completa crea un activo defendible; este es el régimen del ciclo de datos positivo en VRIN de
\cref{sec:data-asset}.

\begin{figure}[htbp]
  \centering
  \includegraphics[width=0.62\linewidth]{ai-build-buy}
  \caption{Marco de adquisición de IA. Eje horizontal: especificidad de la tarea (genérica a
           propietaria). Eje vertical: ventaja en datos propietarios (baja a alta). El cuadrante
           de construcción es el único en el que la propiedad total genera un activo positivo en
           VRIN; los otros tres favorecen comprar u orquestar capacidad externa.}
  \label{fig:ai-build-buy}
\end{figure}

La disciplina financiera es \cref{eq:npv}: construir se justifica solo cuando el valor presente
del beneficio diferenciado supera el valor presente del costo total de propiedad — entrenamiento
inicial, inferencia continua y el costo de mantenimiento que se compone a medida que la
degradación del modelo (riesgo de modelo) requiere reentrenamiento. La opción de comprar tiene un
costo inicial cercano a cero y un precio por unidad completamente conocido; no genera *moat* pero
tampoco deuda de ingeniería.

\begin{example}[Resúmenes de informes con IA: API vs.\ ajuste fino]
La empresa SaaS evalúa dos opciones para una funcionalidad de resumen narrativo generado por IA
que aparece en el informe mensual de análisis de cada cliente.

\textbf{Opción A — Envoltorio de API:} Usar un modelo de lenguaje grande de commodities a través
de API a \money{0.002} por cada mil tokens. Con el volumen típico de informes, esto arroja
aproximadamente \money{10} al mes en costos de inferencia (\money{120} al año). La implementación
requiere dos semanas de ingeniería; no se crea ningún *moat*, porque cualquier competidor puede
llamar al mismo endpoint.

\textbf{Opción B — Ajuste fino sobre datos de uso:} Adaptar un modelo base sobre el corpus de
telemetría multi-tenant propietario de la empresa. Costo único de entrenamiento: \money{15000}.
Costos continuos de inferencia y reentrenamiento: \money{2000} al mes (\money{24000} al año).
El modelo con ajuste fino genera resúmenes que citan los propios datos de referencia y patrones de
flujo de trabajo del cliente — VRIN moderado, porque los datos de entrenamiento superan tres de
las cuatro pruebas de \cref{sec:data-asset}.

El horizonte de decisión es de tres años a la tasa WACC de la empresa del \qty{11}{\percent}. El
costo incremental de la Opción B sobre la Opción A es el costo inicial de \money{15000} más el
diferencial anual de \money{23880} (\money{24000} menos \money{120}), descontado al \qty{11}{\percent}:
\[
  \Delta\mathrm{Cost} \;=\; \money{15000}
    \;+\; \money{23880}\times
    \underbrace{\sum_{t=1}^{3}\frac{1}{(1.11)^{t}}}_{=\;2.4437}
  \;\approx\; \money{73360}.
\]

El beneficio del modelo con ajuste fino fluye a través de la retención: un subconjunto de
aproximadamente 600 clientes activos en informes responde a los resúmenes más ricos, y la
diferenciación reduce el churn de esa cohorte. El NPV es positivo solo si el incremento en CLV
derivado de esta retención supera aproximadamente \money{50} por cliente, porque
\[
  \underbrace{600\;\times\;\money{50}}_{\text{clientes}\;\times\;\Delta\mathrm{CLV}}
  \;\times\; 2.4437
  \;\approx\; \money{73300},
\]
lo que apenas supera el costo incremental. Si la funcionalidad produce solo una retención moderada
— digamos \money{30} por cliente — el NPV es claramente negativo y la Opción A es la correcta.
La implicación: no construya a menos que exista un mecanismo plausible por el cual el modelo con
ajuste fino cree una retención diferenciada \emph{significativa}, no solo una mejora marginal
de calidad.
\end{example}

\begin{admonition}{Nota}
El costo total de propiedad de un sistema de IA personalizado se subestima sistemáticamente en la
etapa de prototipo. El «último cinco por ciento» — endurecimiento de seguridad, SLAs de latencia,
monitoreo de degradación del modelo, pipelines de reentrenamiento y registro de cumplimiento —
representa habitualmente la mayor parte del gasto de ingeniería a lo largo de la vida del sistema.
Un análisis de \cref{eq:npv} que captura solo el costo inicial de entrenamiento siempre favorecerá
la construcción; uno que amortiza la carga total de mantenimiento, típicamente no lo hará.
Modele el costo completo antes de comprometerse.
\end{admonition}

La elección construir vs API es estructuralmente una prueba de *moat*: construya solo cuando el
resultado de la construcción supera el umbral VRIN de \cref{sec:rbv-vrin} y produce un diferencial
sostenido de ROIC\,>\,WACC. En el contexto SaaS, el eje de datos propietarios de
\cref{fig:ai-build-buy} se mapea directamente sobre la lógica del ciclo de datos de
\cref{sec:data-asset}: sin datos propietarios, el ajuste fino compra una mejora temporal de
calidad, no un *moat*. \textcite{agrawal2018prediction} analizan la complementariedad entre la
predicción de bajo costo y el juicio humano escaso como la lógica rectora de la priorización de
inversiones en IA.

\section{Gobernanza de IA y Riesgo de Modelo}\label{sec:ai-governance}
% complexity: 3

Desplegar IA en un contexto comercial introduce una categoría de riesgo que no aparece en un registro
de riesgos de software tradicional. Un error de software convencional produce una salida incorrecta
de una manera conocida y verificable; un modelo de IA generativa produce una salida incorrecta
\emph{con confianza}, en lenguaje natural, de una manera que parece correcta y puede no detectarse
antes de que afecte una decisión. La taxonomía del riesgo tiene cuatro dimensiones principales:
\emph{sesgo} (error sistemático correlacionado con atributos protegidos o sensibles),
\emph{alucinación} (generación segura de contenido factualmente incorrecto), \emph{propiedad
intelectual} (procedencia de los datos de entrenamiento y propiedad de las salidas) y
\emph{regulatorio} (prohibiciones específicas de jurisdicción en casos de uso de IA de alto
riesgo, ejemplificadas por el sistema de clasificación de riesgos de la Ley de IA de la UE).

La pregunta de gobernanza es un problema de agencia en el sentido de \cref{sec:agency-theory}: el
modelo es un agente cuya función objetivo (probabilidad del siguiente token) no está perfectamente
alineada con el objetivo de la empresa (salida precisa y útil). El principal — la empresa que
despliega el modelo — debe diseñar mecanismos de supervisión para detectar el desalineamiento antes
de que cause daño. Esta es precisamente la lógica del sistema de control que \cref{sec:agency-theory}
aplica a los agentes humanos: diseño de incentivos, monitoreo y discrecionalidad acotada. Para los
sistemas de IA, la discrecionalidad acotada se convierte en \emph{humano en el bucle}
(\emph{human-in-the-loop}) para salidas de alto riesgo, y el monitoreo se convierte en validación
automatizada de salidas.

El marco de valor esperado de \cref{sec:expected-value} proporciona la herramienta correcta de
triaje de riesgos. Para cada categoría de salida asistida por IA, la empresa debe estimar la
probabilidad de un error material, el valor en riesgo si ese error se propaga sin detectarse, y el
costo de la capa de revisión humana. Los recursos de gobernanza deben concentrarse donde el
producto de probabilidad y valor en riesgo es mayor.

\begin{example}[Alucinación en proyecciones de ingresos y valor esperado en riesgo]
La empresa SaaS ha desplegado la funcionalidad de resumen de IA de la sección anterior. En un caso,
el modelo emite una proyección de ingresos formulada con seguridad en el informe de un cliente que
sobreestima el crecimiento proyectado del cliente en un \qty{40}{\percent}. El cliente usa la
proyección para planificar plantilla, descubre el error tres meses después y cancela el servicio (churn).

El valor esperado en riesgo de un único incidente de este tipo es el CLV de un cliente que ha
cancelado. Con los parámetros estándar de la empresa (\cref{eq:clv}: \money{252}/año de margen de
contribución, WACC \qty{11}{\percent}, crecimiento \qty{3}{\percent}), CLV\,=\,\money{3150}. Si la
tasa de errores materiales de proyección del modelo se estima en uno por cada 1{,}050 ejecuciones
de informes al mes, y cada error tiene una probabilidad del \qty{50}{\percent} de desencadenar
churn, el costo mensual esperado de churn atribuible a la funcionalidad es \(\money{3150}\times 0.50\approx\money{1575}\).

Un diseño de gobernanza que detecte el error antes de entregarlo al cliente requiere que un analista
revise todos los resúmenes con salidas financieras durante aproximadamente 30 minutos a la semana —
un costo muy inferior a \money{1575} al mes. Más precisamente: el escenario de peor caso de tres
errores no detectados simultáneamente que producen tres clientes cancelados crea un valor esperado
en riesgo de \(3\times\money{3150}=\money{9450}\), que es el piso financiero para la inversión en
gobernanza. La revisión con humano en el bucle de las salidas financieras es claramente de NPV
positivo. Para los resúmenes operativos (adopción de funcionalidades, completaciones de flujos de
trabajo) donde un error es visible para el cliente pero tiene baja consecuencia financiera, los
controles automatizados y la corrección posterior son suficientes.
\end{example}

\begin{admonition}{Advertencia}
«Generado por IA» no es una defensa legal. Si un modelo emite una proyección financiera
materialmente engañosa, una recomendación médica incorrecta o una declaración difamatoria en un
documento de cara al cliente, la responsabilidad recae sobre la empresa que desplegó el modelo, no
sobre el modelo ni su proveedor. Los descargos de responsabilidad contractuales transfieren algo
del riesgo comercial pero no desplazan la responsabilidad legal estatutaria o extracontractual.
Esta es la dimensión regulatoria de la prueba VRIN: un *moat* de datos no sustituible puede ser
simultáneamente una exposición regulatoria no sustituible.
\end{admonition}

El marco de gobernanza se conecta con la atribución en \cref{ch:attribution-and-measurement}: de la
misma manera en que los modelos de atribución requieren un bucle de retroalimentación cerrado entre
la salida del modelo y el resultado real, la gobernanza de IA requiere un mecanismo sistemático para
señalar los errores del modelo, enrutarlos a revisión humana y usar la señal resultante para
reentrenar o restringir el modelo. \textcite{jensen1976agency} establecen la estructura
principal-agente que subyace a estos problemas de control; el contexto de la IA es una nueva
instanciación del desafío de alineación que analizan.

\section{ESG como Problema Estratégico de Datos}\label{sec:esg-data}
% complexity: 3

Los reportes de ESG (factores ambientales, sociales y de gobernanza) suelen enmarcarse como una
obligación de cumplimiento. Su significado estratégico es diferente y mayor: ESG es un \emph{problema
de medición} con la misma estructura que la atribución (\cref{ch:attribution-and-measurement}).
Ambos requieren que la empresa rastree efectos causales de múltiples períodos a través de canales
indirectos, agregue datos de fuentes heterogéneas e informe sobre métricas que afectan las decisiones
de asignación de capital. Ambos están sujetos a mala atribución, selección arbitraria de datos y
incertidumbre del modelo. La disciplina necesaria para medir la huella de carbono con precisión es
exactamente la disciplina necesaria para medir la incrementalidad de marketing con precisión: contra-
factuales claros, condiciones de frontera y auditabilidad.

El marco rector es la \emph{doble materialidad}, codificada en las normas ISSB S1 y S2 y en las
Normas Europeas de Información de Sostenibilidad (ESRS). La doble materialidad exige que una empresa
divulgue temas de sostenibilidad desde dos perspectivas simultáneamente. La \emph{materialidad
financiera} (perspectiva exterior-interior) pregunta cómo los factores ESG — riesgo climático
físico, interrupciones en la cadena de suministro, escasez de talento, endurecimiento regulatorio —
afectan el valor de empresa y los flujos de caja. La \emph{materialidad de impacto} (perspectiva
interior-exterior) pregunta cómo las operaciones de la empresa afectan a las personas y al entorno,
independientemente de si esos impactos actualmente fluyen de vuelta al estado de resultados. Un tema
es reportable si supera el umbral en \emph{cualquiera} de las dos dimensiones; la unión, no la
intersección, determina el perímetro de divulgación.

La lógica del valor para las partes interesadas de \cref{sec:stakeholder-value} proporciona la
justificación estratégica. \textcite{freeman2010strategic} argumenta que la creación de valor
sostenible requiere gestionar simultáneamente las obligaciones con empleados, clientes, proveedores,
comunidades e inversores. La doble materialidad es la arquitectura de reporte que hace esa
obligación medible y auditable.

\begin{example}[Alcances 1/2/3 y la diligencia debida de la Serie~A para la empresa SaaS]
La empresa SaaS con \money{4000000} ARR opera sobre infraestructura en la nube. Las emisiones de
Alcance~1 y Alcance~2 (combustión directa y electricidad adquirida) son casi nulas: la empresa no
posee planta física y su proveedor de nube reporta una cobertura de energía renovable superior al
\qty{90}{\percent}. En un análisis de materialidad puramente financiera, el riesgo climático no es
material en esta etapa: ningún impuesto al carbono ni evento climático físico es plausiblemente
suficientemente grande como para afectar un negocio de \money{4000000} ARR en un horizonte de cinco
años.

Los riesgos ESG materiales son de naturaleza diferente. Las emisiones de Alcance~3 en la cadena de
suministro upstream (fabricación de hardware, carbono incorporado en los centros de datos) son
difíciles de medir pero cada vez más esperadas en la diligencia debida de inversores. En la
dimensión social, la equidad en el talento — equidad salarial por género y etnia, distribución de
opciones sobre acciones — es un riesgo estratégico genuino: con \money{120000} al mes en gasto de
adquisición y un objetivo de 200 clientes al mes, la empresa depende de un equipo pequeño de alta
especialización; las disputas de equity o la rotación crean un riesgo operativo directo. Ninguno de
los dos elementos es financieramente material frente a \money{4000000} ARR de forma aislada, pero
la \emph{materialidad de impacto} requiere divulgación para la diligencia debida de Serie~A en el
marco alineado con ISSB: los inversores institucionales y estratégicos que aplican ISSB S1/S2
esperarán una evaluación de doble materialidad, incluso de una empresa muy por debajo del umbral de
reporte CSRD.

La identidad de beneficio económico de \cref{eq:economic-profit} es el numerador correcto para
evaluar el gasto en ESG: cualquier inversión en ESG que no produzca una reducción en WACC (a través
de menor riesgo regulatorio o menor costo de capital) o un aumento en NOPAT (a través de mejor
retención de talento o menor churn) no crea beneficio económico y debe evaluarse en consecuencia.
Con \money{4000000} ARR, la mayor parte del gasto en ESG cae en el canal del costo de capital —
la reducción en la fricción de financiamiento — más que en el canal de NOPAT.
\end{example}

\begin{admonition}{Nota}
Una puntuación ESG alta no implica ROIC\,>\,WACC. La relación entre el desempeño ESG y el
beneficio económico transcurre a través de dos mecanismos específicos: reducción del costo de
capital (menor prima de riesgo de los inversores que filtran la exposición ESG) y reducción del
riesgo operativo (menor probabilidad de sanciones regulatorias o interrupción del talento). Ningún
mecanismo es automático; ambos son medibles. La advertencia contra confundir las apariencias ante
las partes interesadas con el valor para el accionista se desarrolla completamente en
\cref{ch:value-creation} y no se reitera aquí.
\end{admonition}

El paralelismo con la atribución es más que estructural: las herramientas metodológicas son en gran
medida las mismas. Medir las emisiones de Alcance~3 requiere un modelo de cadena de valor con líneas
de base contrafactuales — ¿cuáles habrían sido las emisiones en ausencia de las operaciones de la
empresa? — que es la misma pregunta que los modelos de atribución formulan sobre el gasto en medios.
Desarrollar capacidad de medición de ESG se compone, por lo tanto, con la capacidad de medición de
atribución; las dos inversiones comparten infraestructura.
\textcite{freeman2010strategic} y \textcite{jensen1976agency} juntos definen la tensión entre
gobernanza y creación de valor que la doble materialidad está diseñada para resolver.
